\documentclass[12pt]{article}
\usepackage[utf8]{inputenc}
\usepackage[T1]{fontenc}
\usepackage{hyperref}
\title{Kursav Language Resources for K-12 Curriculum Development}
\author{Compiled by Claude Sonnet 5.5 (claude-sonnet-5-5) for Andr\'as Kornai}
\date{}
\begin{document}
\maketitle

\section{Introduction}
Kursav (ISO 639-3 \texttt{faj}; also called Faita or Kulsab) is a Madang-family language spoken on the south bank of the upper Ramu River, Usino Bundi district, Madang Province, Papua New Guinea. Search results report only about 10 speakers as of 2015 and classify it as severely endangered; younger people reportedly speak Tok Pisin. The only substantial resources found are archival.

\section{Primary Educational Resources}
No primers, textbooks or school programmes were found.

\section{Linguistic and Reference Materials}
\begin{itemize}
\item Wikipedia overview of Kursav: \url{https://en.wikipedia.org/wiki/Kursav_language}
\item New Guinea World entry for Kursav: \url{https://newguineaworld.linguistik.uzh.ch/families/madang-upper-yuat-river/madang/west-madang/sogeram-and-tomui-rivers/sogeram-river/east-sogeram-river/kursav}
\item Endangered Languages Project entry: \url{https://endangeredlanguages.com/lang/2215?hl=de}
\end{itemize}

\section{Religious and Community Resources}
None found.

\section{Digital and Audio Resources}
Don Daniels's archived Kursav materials are held in the PARADISEC catalogue (search snippets mention a lexicon file and audio recordings; contents were not inspected):
\begin{itemize}
\item \url{https://catalog.paradisec.org.au/repository/DD1/341}
\item \url{https://catalog.paradisec.org.au/collections/DD1/items/136}
\item \url{https://catalog.paradisec.org.au/repository/DD1/377}
\end{itemize}

\section{Curriculum Development Context}
With roughly ten elderly speakers, a school curriculum is not realistic; documentation and revitalisation work based on the PARADISEC archive would be the first step. Access conditions for the archive should be checked with PARADISEC.

\end{document}